\scanpage{126}
\mspage{41}
% archive sheet number: 125

% L01
6. \uline{Application aux faisceaux de cycles évanescents.}

% L02-05
6.1. Soit à nouveau $X$ de type fini sur $S$, \inserted{étant} strictement local, et soit $U$ un ouvert de $X_\eta$, que nous supposons \uline{régulier}. \inserted{Soit $Y=X-U$, et} Supposons qu'il existe un morphisme propre $p:X'\to X$, \deleted{tel que} avec $X'$ régulier, $p^{-1}(Y)$ un diviseur à croisements normaux,

% L05-08 left margin
\diagram{\begin{tikzcd}[ampersand replacement=\&]
X'\arrow[d,no head,"p"']\&U\arrow[l,hook]\arrow[d,no head,"\sim"]\\
X\arrow[d,"f"']\&U\arrow[l,hook]\\
S\&
\end{tikzcd}}

% L06-14
et $U'=p^{-1}(U)\unsure{{-}|\longrightarrow}U$ un isomorphisme. (On sait par Hironaka [\ ] qu'il en est ainsi si $S$ est excellent de car. nulle, et \deleted{sans doute} il semble plausible qu'il en est de même dès que $S$ est excellent, sans restriction de caractéristique. D'ailleurs l'hypothèse ``$S$ excellent'' est anodine, car on s'y ramène pour ttes ces questions qui nous intéressent par passage au complété). \deleted{Sous ces conditions} Choisissons un entier $\ell\ne p$, et considérons, pour \deleted{tout}\inserted{un} entier $q$ \inserted{fixé} les faisceaux

% L15
\[
\boxed{(\underline{\Psi}_\nu^q)_{\mathrm{tame}}=(\underline{\Psi}_\nu^q)^P}
\tag*{6.1.1}
\]

% L16-18
(invariants \inserted{dans $\underline{\Psi}_\nu^q$ sous l'action du} \unsure{par le} p-s-groupe de Sylow $P$ de $I$)\deleted{\unsure{dans}}, \deleted{\unsure{Pour}} relatifs aux coefficients
\[
\Lambda_\nu=\unsure{\mathbb{Z}|Z}/\ell^{\nu+1}\mathbb{Z}\qquad(\nu\geq0).
\]
\inserted{(cf. 2.6 bis)} Ils forment un système projectif de faisceaux de torsion sur $X_0$.

% L19
\uline{Théorème} 6.2 : Avec les notations précédentes :

% L20-23
(i) Les faisceaux $(\underline{\Psi}_\nu^q)_{\mathrm{tame}}$ sont constructibles, et \deleted{il} le système projectif qu'ils forment \deleted{\unsure{$\ell$-adique}} \deleted{\unsure{satisfait} la condition MLAR (SGA 5 V 2.1.1)}, \inserted{est AR-$\ell$-adique (SGA 5 V \unsure{3.7.c|3.2.c})}, donc définissent un faisceau $\ell$-adique constructible (SGA 5 VI 1.1)

% L24
\[
(\underline{\Psi}^q)_{\mathrm{tame}}=\varprojlim_\nu(\underline{\Psi}_\nu^q)_{\mathrm{tame}},
\]
